\documentclass[11pt, a4paper]{article}
% Layout / typography / colour
\usepackage[a4paper,margin=2cm]{geometry} % Page size and margins
\usepackage{graphicx}                    % Include images
\usepackage{float}                       % Improved float placement control
\usepackage{array}                       % Extended table column definitions
\usepackage{multirow}                    % Multi-row cells in tables
\usepackage{subcaption}                  % Subfigures with captions
\usepackage{bm}                          % Bold maths symbols
\usepackage{microtype}                   % Improved justification and kerning
\usepackage{xcolor}                      % Colour support
% Maths
\usepackage{amsmath}                     % Core maths environments
\usepackage{amsfonts}                    % Additional maths fonts
\usepackage{amssymb}                     % Extended maths symbols
\usepackage{xfp}                         % Automatic maths calculations
\usepackage{mathtools}                   % Extensions and fixes for amsmath
\usepackage{siunitx}                     % Consistent units and quantities
% Environments / code
\usepackage{amsthm}                      % Theorem-like environments
\usepackage{thmtools}                    % Enhanced theorem management
\usepackage{listings}                    % Source code listings
% Diagrams
\usepackage{tikz}                        % Vector graphics and diagrams
\usepackage{circuitikz}                  % Electrical circuit diagrams
\usepackage{pgfplots}                    % Plots and data visualisation
\pgfplotsset{compat=1.18}                % PGFPlots compatibility level
\usetikzlibrary{arrows.meta,positioning,calc,shapes.geometric} % Common TikZ helpers
% References / links
\usepackage{cite}                        % Compressed numeric citations
\usepackage[colorlinks=true,linkcolor=black,citecolor=black,urlcolor=black]{hyperref} % Hyperlinks in PDF
% Glossaries
\usepackage{cleveref}                    % Intelligent cross-references
% Structure
\usepackage{fancyhdr}                    % Custom headers and footers
\usepackage{sectsty}                     % Section heading styling
\usepackage{titling}                     % Control over title formatting
% Lists / utilities
\usepackage{enumitem}                    % Customisable list environments
\usepackage{xparse}                      % Modern command definitions
\usepackage{etoolbox}                    % Programming utilities and conditionals

% Listings base configuration. Listing Usage: \begin{lstlisting}[style=style, caption={Caption}, label={lst:label}]\end{lstlisting}
\lstset{
  basicstyle=\ttfamily\small,
  keywordstyle=\color{blue}\bfseries,
  stringstyle=\color{orange},
  commentstyle=\color{gray}\itshape,
  numberstyle=\tiny\color{gray},
  numbers=left,
  numbersep=8pt,
  frame=single,
  rulecolor=\color{black},
  breaklines=true,
  breakatwhitespace=true,
  tabsize=2,
  showstringspaces=false,
  captionpos=b,
  keepspaces=true
}
% C
\lstdefinestyle{cStyle}{
  language=C,
  morekeywords={uint8_t,uint16_t,uint32_t,size_t}
}
% C++
\lstdefinestyle{cppStyle}{
  language=C++,
  morekeywords={constexpr,nullptr,override,noexcept}
}
% Python
\lstdefinestyle{pythonStyle}{
  language=Python,
  morekeywords={self}
}
% MATLAB
\lstdefinestyle{matlabStyle}{
  language=Matlab,
  morekeywords={end,function}
}
% Java
\lstdefinestyle{javaStyle}{
  language=Java
}
% Bash
\lstdefinestyle{bashStyle}{
  language=bash,
  morekeywords={sudo,cd,ls,grep,awk,sed}
}
% LaTeX
\lstdefinestyle{latexStyle}{
  language={[LaTeX]TeX},
  morekeywords={usepackage,newcommand,begin,end}
}
% Inline code
\newcommand{\code}[1]{\lstinline{#1}}

% Dark Blue heading colour
\definecolor{header}{RGB}{30,50,100}
\sectionfont{\color{header}}
\subsectionfont{\color{header}}
\subsubsectionfont{\color{header}}

\title{Project SunSat: Project Enki \\ \textbf{Critical Design Review}}
\author{R.Khawaja, L.Peintre, M.Tsui}
\date{}

\begin{document}

\bibliographystyle{IEEEtran}
\maketitle

\begin{abstract}
Second year computer system / robotics engineers, Rohaan Khawaj, Marcus Tsui \& Leo Peintre, propose the development of a CubeSat payload designed for the acquisition of magnetic field data for regional crustal mapping. The authors aim to build upon prior experience as avionics engineers in Project SunSat. The project aims to design, prototype, and evaluate a payload capable of collecting low-noise, geomagnetic field data from stratospheric altitudes (30-35km) via a weather balloon. The collected data will be suitable for post processing and geological analysis. The scope of this project includes onboard data handling, telemetry to a ground control station, and data processing for use in geological studies. The primary outcome of this initial proposal is to communicate our project to TBD.

Skip to content
logo
niri
Input
Type to start searching
GitHub

    v26.04
    26k
    975

    Usage
    Configuration
        Introduction
        Input
        Outputs
        Key Bindings
        Switch Events
        Layout
        Named Workspaces
        Miscellaneous
        Window Rules
        Layer Rules
        Animations
        Gestures
        Recent Windows
        Debug Options
        Include
    Development

    Overview
    Keyboard
        Layout
        Repeat
        Num Lock
    Pointing Devices
    General Settings
        disable-power-key-handling
        warp-mouse-to-focus
        focus-follows-mouse
        workspace-auto-back-and-forth
        mod-key, mod-key-nested

Input
Overview#

In this section you can configure input devices like keyboard and mouse, and some input-related options.

There's a section for each device type: keyboard, touchpad, mouse, trackpoint, trackball, tablet, touch. Settings in those sections will apply to every device of that type. Currently, there's no way to configure specific devices individually (but that is planned).

All settings at a glance:

input {
    keyboard {
        xkb {
            // layout "us"
            // variant "colemak_dh_ortho"
            // options "compose:ralt,ctrl:nocaps"
            // model ""
            // rules ""
            // file "~/.config/keymap.xkb"
        }

        // repeat-delay 600
        // repeat-rate 25
        // track-layout "global"
        numlock
    }

    touchpad {
        // off
        tap
        // dwt
        // dwtp
        // drag false
        // drag-lock
        natural-scroll
        // accel-speed 0.2
        // accel-profile "flat"
        // scroll-factor 1.0
        // scroll-factor vertical=1.0 horizontal=-2.0
        // scroll-method "two-finger"
        // scroll-button 273
        // scroll-button-lock
        // tap-button-map "left-middle-right"
        // click-method "clickfinger"
        // left-handed
        // disabled-on-external-mouse
        // middle-emulation
    }

    mouse {
        // off
        // natural-scroll
        // accel-speed 0.2
        // accel-profile "flat"
        // scroll-factor 1.0
        // scroll-factor vertical=1.0 horizontal=-2.0
        // scroll-method "no-scroll"
        // scroll-button 273
        // scroll-button-lock
        // left-handed
        // middle-emulation
    }

    trackpoint {
        // off
        // natural-scroll
        // accel-speed 0.2
        // accel-profile "flat"
        // scroll-method "on-button-down"
        // scroll-button 273
        // scroll-button-lock
        // left-handed
        // middle-emulation
    }

    trackball {
        // off
        // natural-scroll
        // accel-speed 0.2
        // accel-profile "flat"
        // scroll-method "on-button-down"
        // scroll-button 273
        // scroll-button-lock
        // left-handed
        // middle-emulation
    }

    tablet {
        // off
        map-to-output "eDP-1"
        // map-to-focused-output
        // map-to-focused-window
        // left-handed
        // calibration-matrix 1.0 0.0 0.0 0.0 1.0 0.0
    }

    touch {
        // off
        map-to-output "eDP-1"
        // calibration-matrix 1.0 0.0 0.0 0.0 1.0 0.0
    }

    // disable-power-key-handling
    // warp-mouse-to-focus
    // focus-follows-mouse max-scroll-amount="0%"
    // workspace-auto-back-and-forth

    // mod-key "Super"
    // mod-key-nested "Alt"
}

Keyboard#
Layout#

In the xkb section, you can set layout, variant, options, model and rules. These are passed directly to libxkbcommon, which is also used by most other Wayland compositors. See the xkeyboard-config(7) manual for more information.

input {
    keyboard {
        xkb {
            layout "us"
            variant "colemak_dh_ortho"
            options "compose:ralt,ctrl:nocaps"
        }
    }
}

Tip

Since: 25.02

Alternatively, you can directly set a path to a .xkb file containing an xkb keymap. This overrides all other xkb settings.

input {
    keyboard {
        xkb {
            file "~/.config/keymap.xkb"
        }
    }
}

Note

Since: 25.08

If the xkb section is empty (like it is by default), niri will fetch xkb settings from systemd-localed at org.freedesktop.locale1 over D-Bus. This way, for example, system installers can dynamically set the niri keyboard layout. You can see this layout in localectl and change it with localectl set-x11-keymap, for example:

$
 localectl set-x11-keymap "us" "" "colemak_dh_ortho" "compose:ralt,ctrl:nocaps"
$
 localectl
System
 Locale: LANG=en_US.UTF-8
               LC_NUMERIC=ru_RU.UTF-8
               LC_TIME=ru_RU.UTF-8
               LC_MONETARY=ru_RU.UTF-8
               LC_PAPER=ru_RU.UTF-8
               LC_MEASUREMENT=ru_RU.UTF-8
    VC Keymap: us-colemak_dh_ortho
   X11 Layout: us
  X11 Variant: colemak_dh_ortho
  X11 Options: compose:ralt,ctrl:nocaps

By default, localectl will set the TTY keymap to the closest match of the XKB keymap. You can prevent that with a --no-convert flag, for example: localectl set-x11-keymap --no-convert "us,ru".

These settings are picked up by some other programs too, like GDM.

When using multiple layouts, niri can remember the current layout globally (the default) or per-window. You can control this with the track-layout option.

    global: layout change is global for all windows.
    window: layout is tracked for each window individually.

input {
    keyboard {
        track-layout "global"
    }
}

Repeat#

Delay is in milliseconds before the keyboard repeat starts. Rate is in characters per second.

input {
    keyboard {
        repeat-delay 600
        repeat-rate 25
    }
}

Num Lock#

Since: 25.05

Set the numlock flag to turn on Num Lock automatically at startup.

You might want to disable (comment out) numlock if you're using a laptop with a keyboard that overlays Num Lock keys on top of regular keys.

input {
    keyboard {
        numlock
    }
}

Pointing Devices#

Most settings for the pointing devices are passed directly to libinput. Other Wayland compositors also use libinput, so it's likely you will find the same settings there. For flags like tap, omit them or comment them out to disable the setting.

A few settings are common between input devices:

    off: if set, no events will be sent from this device.

A few settings are common between touchpad, mouse, trackpoint, and trackball:

    natural-scroll: if set, inverts the scrolling direction.
    accel-speed: pointer acceleration speed, valid values are from -1.0 to 1.0 where the default is 0.0.
    accel-profile: can be adaptive (the default) or flat (disables pointer acceleration).
    scroll-method: when to generate scroll events instead of pointer motion events, can be no-scroll, two-finger, edge, or on-button-down. The default and supported methods vary depending on the device type.
    scroll-button: Since: 0.1.10 the button code used for the on-button-down scroll method. You can find it in libinput debug-events.
    scroll-button-lock: Since: 25.08 when enabled, the button does not need to be held down. Pressing once engages scrolling, pressing a second time disengages it, and double click acts as single click of the the underlying button.
    left-handed: if set, changes the device to left-handed mode.
    middle-emulation: emulate a middle mouse click by pressing left and right mouse buttons at once.

Settings specific to touchpads:

    tap: tap-to-click.
    dwt: disable-when-typing.
    dwtp: disable-when-trackpointing.
    drag: Since: 25.05 can be true or false, controls if tap-and-drag is enabled.
    drag-lock: Since: 25.02 if set, lifting the finger off for a short time while dragging will not drop the dragged item. See the libinput documentation.
    tap-button-map: can be left-right-middle or left-middle-right, controls which button corresponds to a two-finger tap and a three-finger tap.
    click-method: can be button-areas or clickfinger, changes the click method.
    disabled-on-external-mouse: do not send events while external pointer device is plugged in.

Settings specific to touchpad and mouse:

    scroll-factor: Since: 0.1.10 scales the scrolling speed by this value.

    Since: 25.08 You can also override horizontal and vertical scroll factor separately like so: scroll-factor horizontal=2.0 vertical=-1.0

Settings specific to tablet and touch:

    calibration-matrix: set to six floating point numbers to change the calibration matrix. See the LIBINPUT_CALIBRATION_MATRIX documentation for examples.
        Since: 25.02 for tablet
        Since: 25.11 for touch

Tablets and touchscreens are absolute pointing devices that can be mapped to a specific output like so:

input {
    tablet {
        map-to-output "eDP-1"
    }

    touch {
        map-to-output "eDP-1"
    }
}

Valid output names are the same as the ones used for output configuration.

Since: 0.1.7 When a tablet is not mapped to any output, it will map to the union of all connected outputs, without aspect ratio correction.

Settings specific to tablet:

    map-to-focused-output: Since: 26.04 will map the tablet to the focused output, takes precedence over map-to-output.

    map-to-focused-window: Since: next release will map the tablet to the focused window's geometry, takes precedence over map-to-focused-output and map-to-output. Falls back to those when no window is focused (for example, in the overview).

    When the tablet is also mapped to a specific output via map-to-output, the map-to-focused-window flag will map the tablet to the active window on that output. If the tablet isn't mapped to any specific output, it will map the tablet to the current focused window regardless of where it is.

General Settings#

These settings are not specific to a particular input device.
disable-power-key-handling#

By default, niri will take over the power button to make it sleep instead of power off. Set this if you would like to configure the power button elsewhere (i.e. logind.conf).

input {
    disable-power-key-handling
}

warp-mouse-to-focus#

Makes the mouse warp to newly focused windows.

Does not make the cursor visible if it had been hidden.

input {
    warp-mouse-to-focus
}

By default, the cursor warps separately horizontally and vertically. I.e. if moving the mouse only horizontally is enough to put it inside the newly focused window, then the mouse will move only horizontally, and not vertically.

Since: 25.05 You can customize this with the mode property.

    mode="center-xy": warps by both X and Y coordinates together. So if the mouse was anywhere outside the newly focused window, it will warp to the center of the window.
    mode="center-xy-always": warps by both X and Y coordinates together, even if the mouse was already somewhere inside the newly focused window.

input {
    warp-mouse-to-focus mode="center-xy"
}

focus-follows-mouse#

Focuses windows and outputs automatically when moving the mouse over them.

input {
    focus-follows-mouse
}

Since: 0.1.8 You can optionally set max-scroll-amount. Then, focus-follows-mouse won't focus a window if it will result in the view scrolling more than the set amount. The value is a percentage of the working area width.

input {
    // Allow focus-follows-mouse when it results in scrolling at most 10% of the screen.
    focus-follows-mouse max-scroll-amount="10%"
}

input {
    // Allow focus-follows-mouse only when it will not scroll the view.
    focus-follows-mouse max-scroll-amount="0%"
}

workspace-auto-back-and-forth#

Normally, switching to the same workspace by index twice will do nothing (since you're already on that workspace). If this flag is enabled, switching to the same workspace by index twice will switch back to the previous workspace.

Niri will correctly switch to the workspace you came from, even if workspaces were reordered in the meantime.

input {
    workspace-auto-back-and-forth
}

mod-key, mod-key-nested#

Since: 25.05

Customize the Mod key for key bindings. Only valid modifiers are allowed, e.g. Super, Alt, Mod3, Mod5, Ctrl, Shift.

By default, Mod is equal to Super when running niri on a TTY, and to Alt when running niri as a nested winit window.

Note

There are a lot of default bindings with Mod, none of them "make it through" to the underlying window. You probably don't want to set mod-key to Ctrl or Shift, since Ctrl is commonly used for app hotkeys, and Shift is used for, well, regular typing.

// Switch the mod keys around: use Alt normally, and Super inside a nested window.
input {
    mod-key "Alt"
    mod-key-nested "Super"
}

Made with Material for MkDocs
\end{abstract}

\tableofcontents
\pagebreak

% Header Style
\pagestyle{fancy}
\fancyhf{}
\fancyhead[L]{Project SunSat: Project Enki}
\fancyhead[R]{Critical Design Review}
\fancyfoot[C]{\thepage}

\section{Introduction}
Project Enki leverages the CubeSat platform in order to deliver an affordable payload that can provide accurate data for further scientific analysis despite using less resources. The CubeSat platform allows for rapid development, low costs, and substantially cheaper launches. These cheaper launches can be achieved by using the CubeSat as a secondary payload and are further helped by its relatively low mass. This is further  
\section{Component List}

\section{Data Use Case \& Pipeline}

\section{Electrical}
\subsection{Electrical Power System}
\subsection{Circuit Schematic}
\subsection{PCB Design}
\subsection{Thermal Considerations}
\subsection{Shielding \& Interference Prevention}

\section{Command \& Data Handling}
\subsection{}

\section{Telemetry}

\section{Attitude Determination \& Orbital Control System}

\section{Guidance, Navigation \& Control}

\section{Manufacturing Technique}

\section{Timeline \& Budget Considerations}

\pagebreak
\bibliography{references}
\end{document}
